\scanpage{225}
\mspage{80}

\deleted{\unsure{6)}} Tous ces exemples sont de $1^{\mathrm{ere}}$ espèce, les ensembles ordonnés qui interviennent sont de la forme $\unsure{\mathfrak P|\mathcal P}(P_0),\mathfrak P(Q_0)$, et la relation entre eux est déduite d'une relation $R_0$ entre $P_0$ et $Q_0$. \unsure{voici|Voici} maintenant des exemples plus complexes.

6) Soit $M$ une catégorie (``grosse'', le plus souvent \ldots), On prend $\unsure{\mathfrak P|P}$

\[
 \unsure{\mathfrak P|P}=Q=\mathfrak P(\mathrm{Fl}(M)).
\]

La relation $R\subset P\times Q$ est la relation de Quillen, que j'écris, pour $\Phi,\Psi\in\mathfrak P$

\[
 \begin{aligned}
 &\deleted{\unsure{B\subset A,}}\\
 &\deleted{\unsure{\Phi}}\Psi\subset\Phi_*\,\unsure{,|?}\quad\deleted{\unsure{Q|R}}
 \end{aligned}
\]

ou encore de façon équivalente

\[
 \Phi\subset\Psi^*
\]

($\unsure{\Psi|\Phi|\Psi^{+}}$ a la LLP pour $\Psi$, $\Phi$ la RLP pour $\Psi$). \deleted{\unsure{Les}} \deleted{Donc} les applications $\alpha_0,\beta_0$ sont données par

\[
 \beta_0(\Phi)=\Phi_*,\qquad \alpha_0(\Psi)=\Psi^*.
\]

L'idée ici ne viendrait à personne d'introduire $\mathfrak P(P),\mathfrak P(Q)$ \deleted{\unsure{et}} c'est déjà assez gros comme ça ! La notation pour $\Phi^{\unsure{\sharp}}$ est ici $\unsure{\widetilde\Phi|\overline\Phi}$, pour $\unsure{\widehat{\Psi}|\widetilde{\Psi}|\Psi}^\flat$ c'est $\overline\Psi$. Les parties closes sont celles que j'appelle $Q$-closes à g. et à droite.

La relation $R$, que j'écris

\[
 \Phi\longleftrightarrow\Psi,
\]
% archive sheet number: 225
